\section{Conclusion}
\label{Conclusion}
This thesis evaluated a barrier control algorithm against an FAS controller on a closed-loop figure-eight track across four traffic densities under varying initial conditions, using CARLA, a physics based simulator. Both evaluated controllers were subjected to identical simulation conditions, ensuring that differences in controller performance in terms of KPIs, could be solely attributed to the controllers. The results show that the barrier controller outperforms the FAS controller in throughput, delay, and energy consumption, across all traffic densities. The results can be attributed to the formed steady-state platoon, which allows continuous crossing of the intersection, eliminating idle time and stop-and-go behavior, which are in fact present during simulations of the FAS controller. The energy consumption reduction highlighted that stop-and-go behavior at intersections is a driver of increased energy consumption of electric vehicles. 
These performance gains, however, come at a safety cost at higher traffic densities, making safety the only KPI where the FAS controller performs better. Nevertheless, the violations only occurred during the transient phase, which is the phase before the steady-state platoon is formed. After the platoon was formed, the barrier showed no violations, indicating that the barrier controller is capable of safe operation, but needs parameter tuning for higher traffic densities. 
\\\\
Future research could improve and extend this project in five directions. First, research should focus on the adjustment of the barrier controller parameters to eliminate safety violations. This tuning can be in terms of either more aggressive fixed values or as an adaptive function based on traffic density. Second, future research should extend to a double-intersection, to assess if the barrier controller is scalable beyond the single-intersection. Third, the performance of the barrier controller should be evaluated under sensor noise, since the simulation assumes perfect distance estimation, whereas sensors are subject to measurement noise, which may affect safety and overall performance of the barrier controller. Fourth, for actual energy performance the energy model should be made more accurate by removing the stated assumptions. 
Fifth, the barrier control algorithm should be evaluated in a controlled laboratory or in the real world. 
